\documentclass[12pt, a4paper]{report}

\usepackage[utf8]{inputenc}
\usepackage[spanish]{babel}
\usepackage{graphicx}
\usepackage{geometry}
\usepackage{hyperref}
\usepackage{listings}
\usepackage{color}
\usepackage{amsmath}
\usepackage{booktabs}
\usepackage{longtable}
\usepackage{float}

\geometry{left=25mm, right=25mm, top=25mm, bottom=25mm}

\definecolor{codegreen}{rgb}{0,0.6,0}
\definecolor{codegray}{rgb}{0.5,0.5,0.5}
\definecolor{codepurple}{rgb}{0.58,0,0.82}
\definecolor{backcolour}{rgb}{0.95,0.95,0.92}

\lstdefinestyle{mystyle}{
    backgroundcolor=\color{backcolour},
    commentstyle=\color{codegreen},
    keywordstyle=\color{magenta},
    numberstyle=\tiny\color{codegray},
    stringstyle=\color{codepurple},
    basicstyle=\ttfamily\footnotesize,
    breakatwhitespace=false,
    breaklines=true,
    captionpos=b,
    keepspaces=true,
    numbers=left,
    numbersep=5pt,
    showspaces=false,
    showstringspaces=false,
    showtabs=false,
    tabsize=2
}
\lstset{style=mystyle}

\title{\textbf{NEVEN: Sistema Políglota de Cómputo Científico\\Integrado en Microsoft Excel\\con R, Julia, Python, NEVEN Studio y Motor RAG}}
\author{Minor Bonilla G.\\BukloLAB\\Tesis de Maestría -- Universidad de Costa Rica}
\date{Septiembre de 2026}

\begin{document}

\maketitle

\begin{abstract}
NEVEN es un sistema de cómputo científico políglota que integra R 4.4.1, Julia 1.12.6 y Python 3.12 como motores de scripting embebidos dentro de Microsoft Excel, utilizando una arquitectura de cuatro capas con procesos aislados comunicados mediante Named Pipes y Protocol Buffers. El sistema expone más de 200 funciones definidas por el usuario (UDFs) accesibles directamente desde celdas de Excel, un visor HTML embebido basado en WebView2 para visualizaciones interactivas, notebooks Pluto.jl reactivos, un motor RAG (Retrieval-Augmented Generation) con más de 300 entidades ontológicas, y NEVEN Studio: un modo standalone con siete tabs especializados (SQL, Data Studio, Run Script, Data Lab, Presentaciones, IA, Ayuda). La validación se sustenta en una suite de 342 pruebas automatizadas con 100\% de tests pasando, incluyendo pruebas basadas en propiedades. El sistema obtuvo una calificación de calidad de \textbf{9.71/10} en la evaluación doctoral, demostrando madurez en arquitectura, calidad de código, y completitud funcional.
\end{abstract}

\tableofcontents
\newpage

% ═══════════════════════════════════════════════════════════════════════════
\chapter{Introducción}
% ═══════════════════════════════════════════════════════════════════════════

\section{Contexto y Motivación}

En las ciencias de datos modernas existe una brecha entre el poder del procesamiento estadístico avanzado (R, Julia, Python) y la accesibilidad en la capa de negocio (Microsoft Excel). Las soluciones existentes como BERT (Basic Excel R Toolkit) adolecen de problemas de estabilidad al inyectar librerías dinámicas directamente en el proceso de Excel, donde una violación de segmento del motor analítico provoca el colapso del proceso anfitrión.

NEVEN resuelve este problema mediante una arquitectura de cuatro capas con aislamiento de procesos: el XLL (add-in de Excel) se comunica con procesos hijo independientes (ControlR.exe, ControlJulia.exe, ControlPython.exe) a través de Named Pipes, garantizando que un fallo en cualquier motor no afecte la estabilidad de Excel.

Más allá del aislamiento técnico, NEVEN introduce un paradigma de accesibilidad: el usuario de negocio puede ejecutar análisis estadísticos avanzados escribiendo fórmulas familiares en Excel (\texttt{=NEVEN.r("regresion\_lineal", datos, "html")}), mientras el usuario técnico dispone de NEVEN Studio para scripting libre, consultas SQL sobre datos de Excel, y un asistente de IA contextual.

\section{Objetivos}

\begin{enumerate}
    \item Aislar la ejecución del código analítico en procesos separados del entorno UI de Excel.
    \item Integrar tres motores de lenguaje (R, Julia, Python) con una interfaz unificada.
    \item Exponer más de 200 funciones estadísticas, matemáticas y de machine learning como UDFs de Excel.
    \item Proporcionar NEVEN Studio como modo standalone con tabs especializados para usuarios técnicos.
    \item Implementar un motor RAG con ontología de dominio para asistencia contextual.
    \item Integrar visualizaciones interactivas mediante WebView2, Plotly, y notebooks Pluto.jl reactivos.
    \item Garantizar la reproducibilidad de análisis mediante exportación a notebooks y Quarto.
    \item Mantener latencia sub-milisegundo ($<$1ms) en la comunicación IPC.
\end{enumerate}

\section{Contribuciones}

Las contribuciones principales de este trabajo son:
\begin{itemize}
    \item Una arquitectura IPC de cuatro capas basada en Named Pipes y Protocol Buffers que soporta tres lenguajes de programación con latencia $<$1ms.
    \item La corrección de bugs críticos en la C API de Julia 1.12 (\texttt{jl\_arrayset} y \texttt{jl\_array\_ptr} para tipos primitivos inline).
    \item Un sandbox de seguridad que bloquea más de 90 patrones de ejecución peligrosa en R, Julia y Python.
    \item NEVEN Studio: un modo standalone con siete tabs especializados (SQL, Data Studio, Run Script, Data Lab, Presentaciones, IA, Ayuda).
    \item Un motor RAG con más de 300 entidades ontológicas para asistencia contextual y generación de código.
    \item Data Lab: pipeline de descubrimiento de datos con DuckDB embebido y motor de preguntas en lenguaje natural.
    \item Documentación completa de 17 capítulos en formato Docusaurus.
\end{itemize}

% ═══════════════════════════════════════════════════════════════════════════
\chapter{Arquitectura del Sistema}
% ═══════════════════════════════════════════════════════════════════════════

\section{Estructura de Cuatro Capas}

NEVEN implementa una topología de cuatro capas:

\begin{enumerate}
    \item \textbf{Capa de Presentación (XLL + Studio):} Add-in de Excel compilado como DLL en C++17 con MSVC 2022 (NEVEN64.xll). Registra funciones UDF, gestiona el visor WebView2, y coordina la comunicación con los motores. NEVEN Studio extiende esta capa con un Office Add-in (HTML/JS) para el modo standalone.
    
    \item \textbf{Capa de Orquestación:} El singleton \texttt{RJ2XCL\_Engine} coordina el ciclo de vida de los motores, el LanguageManager selecciona el motor apropiado según el prefijo de función (R., J., P.), y el servidor HTTP Python (puerto 5555) maneja las peticiones de NEVEN Studio.
    
    \item \textbf{Capa de Transporte (IPC):} Named Pipes de Windows con serialización Protocol Buffers v21.12. Comunicación síncrona bidireccional con timeouts configurables por lenguaje. Latencia medida: $<$1ms.
    
    \item \textbf{Capa de Motores:} Procesos hijo independientes --- ControlR.exe (R 4.4.1), ControlJulia.exe (Julia 1.12.6 con sysimage precompilada de 415MB), y ControlPython.exe (Python 3.12 via Stable ABI) --- que ejecutan código y retornan resultados serializados.
\end{enumerate}

\section{Diagrama de Componentes}

\begin{figure}[H]
\centering
\begin{verbatim}
┌─────────────────────────────────────────────────────────────────┐
│                    Microsoft Excel (Host)                       │
├─────────────────────────────────────────────────────────────────┤
│  NEVEN64.xll          │  NEVENRibbon.dll  │  NEVEN Studio      │
│  (C++17, MSVC 2022)   │  (COM Add-in)     │  (Office Add-in)   │
├───────────────────────┴───────────────────┴─────────────────────┤
│                    RJ2XCL_Engine (Singleton)                    │
│  ┌─────────────┐  ┌─────────────┐  ┌─────────────────────────┐  │
│  │LanguageMgr  │  │ ViewerMgr   │  │ HTTP Server (5555)      │  │
│  │ (R/J/P)     │  │ (WebView2)  │  │ (Python + DuckDB)       │  │
│  └──────┬──────┘  └─────────────┘  └─────────────────────────┘  │
├─────────┼───────────────────────────────────────────────────────┤
│         │           Named Pipes + Protocol Buffers              │
├─────────┼───────────────────────────────────────────────────────┤
│    ┌────┴────┐    ┌────────────┐    ┌─────────────┐             │
│    │ControlR │    │ControlJulia│    │ControlPython│             │
│    │ (R 4.4) │    │ (Jl 1.12.6)│    │ (Py 3.12)   │             │
│    └─────────┘    └────────────┘    └─────────────┘             │
└─────────────────────────────────────────────────────────────────┘
\end{verbatim}
\caption{Arquitectura de cuatro capas de NEVEN}
\end{figure}

\section{Componentes Principales}

\begin{table}[H]
\centering
\begin{tabular}{lll}
\toprule
\textbf{Componente} & \textbf{Tecnología} & \textbf{Función} \\
\midrule
NEVEN64.xll & C++17, MSVC 2022 & Add-in de Excel, registro de UDFs \\
ControlR.exe & C++ + R 4.4.1 & Motor R con embedding via R API \\
ControlJulia.exe & C++ + Julia 1.12.6 & Motor Julia + sysimage 415MB \\
ControlPython.exe & C++ + Python 3.12 & Motor Python via Stable ABI \\
neven\_http\_server.py & Python + FastAPI & Servidor HTTP para Studio \\
ViewerManager & C++ + WebView2 SDK & Visor HTML embebido \\
PlutoManager & C++ & Gestión del servidor Pluto.jl \\
PresentationBuilder & C++ & Generador reveal.js \\
RAG Engine & Python + FAISS & Motor de recuperación semántica \\
DuckDB & SQL embebido & Motor SQL para Data Lab \\
SandboxVerifier & C++ & Validación de seguridad de código \\
ConfigService & C++ + JSON & Configuración centralizada \\
\bottomrule
\end{tabular}
\caption{Componentes principales del sistema NEVEN}
\end{table}

\section{Protocol Buffers para IPC}

El intercambio de datos utiliza Protocol Buffers con un esquema extensible:

\begin{lstlisting}[language=C++, caption={Esquema Protobuf para variable.proto}]
message Variable {
  oneof value {
    int32 integer = 5;
    double real = 6;
    string str = 7;
    bool boolean = 8;
    Array arr = 10;
    HtmlContent html_content = 16;
    PlotlyJson plotly_json = 17;
  }
  string name = 15;
  string label = 18;
  string type_hint = 19;
}
\end{lstlisting}

Los campos \texttt{HtmlContent} (field 16) y \texttt{PlotlyJson} (field 17) fueron agregados para transportar visualizaciones interactivas al visor WebView2.

\section{Gestión de Memoria: Patrón RAII}

La interfaz \texttt{XLOPER12} de Excel requiere gestión manual de memoria. NEVEN encapsula toda interacción en la clase \texttt{RaiiXlOper}, garantizando liberación determinista:

\begin{lstlisting}[language=C++, caption={Patrón RAII para XLOPER12}]
{
    RaiiXlOper result;
    Excel12(xlUDF, &result, 1, arg1);
    // Al salir del bloque, el destructor libera
    // la memoria de XLOPER12 automaticamente.
}
\end{lstlisting}

% ═══════════════════════════════════════════════════════════════════════════
\chapter{Funcionalidades Implementadas}
% ═══════════════════════════════════════════════════════════════════════════

\section{Funciones Estadísticas en R}

Se implementaron 32 archivos R con más de 90 procedimientos estadísticos, accesibles como \texttt{=NEVEN.r("funcion", datos, "TipoOutput")}:

\begin{table}[H]
\centering
\begin{tabular}{llc}
\toprule
\textbf{Módulo} & \textbf{Paquete R} & \textbf{Procedimientos} \\
\midrule
Regresión Lineal & stats, stargazer & 12 \\
Regresión Binaria & stats & 8 \\
Datos de Panel & plm & 15 \\
Series de Tiempo & tseries, mFilter & 12 \\
Pronóstico & forecast, auto.arima & 10 \\
Efectos Mixtos & lme4 & 8 \\
Supervivencia & survival & 7 \\
Psicometría & psych & 7 \\
Bayesiana & rstanarm & 8 \\
Estadística Descriptiva & Hmisc & 12 \\
Supuestos & car, lmtest & 12 \\
Visualización & ggplot2, plotly & 8 \\
\bottomrule
\end{tabular}
\caption{Módulos estadísticos R implementados}
\end{table}

\section{Funciones Científicas en Julia}

Se implementaron 4 módulos Julia (J4XCL) con 61 procedimientos, accesibles como \texttt{=NEVEN.j("funcion", datos, "TipoOutput")}:

\begin{table}[H]
\centering
\begin{tabular}{llc}
\toprule
\textbf{Módulo} & \textbf{Área} & \textbf{Procedimientos} \\
\midrule
J4XCL\_Math & Álgebra Lineal, Cálculo, EDO & 23 \\
J4XCL\_ML & Machine Learning, Clustering & 18 \\
J4XCL\_Optimize & Optimización, LP, NLP & 12 \\
J4XCL\_Utils & Transformación, Utilidades & 8 \\
\bottomrule
\end{tabular}
\caption{Módulos Julia (J4XCL) implementados}
\end{table}

\textbf{Nota sobre Julia sysimage:} Para eliminar el cold start de Julia (que originalmente tomaba minutos), NEVEN incluye una sysimage precompilada (\texttt{neven\_julia.dll}, 415MB) que reduce el tiempo de inicio a segundos.

\section{Funciones Python}

Python está integrado como tercer motor via Stable ABI (\texttt{python3.dll}), permitiendo compatibilidad con múltiples versiones de Python sin recompilación. El motor expone funciones accesibles como \texttt{=NEVEN.py("funcion", datos, "TipoOutput")}.

\section{Visor WebView2 Embebido}

El sistema integra Microsoft Edge WebView2 como visor HTML embebido dentro de Excel, permitiendo renderizar visualizaciones interactivas (Plotly, D3.js, htmlwidgets) sin abrir un navegador externo. El visor soporta:

\begin{itemize}
    \item Ventanas flotantes modeless asociadas a la ventana de Excel.
    \item Comunicación bidireccional JavaScript$\leftrightarrow$C++ via PostMessage.
    \item Sandbox de seguridad con navegación restringida.
    \item Múltiples instancias simultáneas con evicción FIFO.
    \item Integración con NEVEN Studio para visualización de resultados.
\end{itemize}

\section{Notebooks Pluto.jl Reactivos}

Para usuarios técnicos, el sistema ofrece notebooks Pluto.jl reactivos:

\begin{itemize}
    \item Servidor Pluto.jl lanzado como proceso Julia independiente.
    \item 13 notebooks precargados (7 R via RCall.jl, 5 Julia nativo, 1 mixto).
    \item Pipeline de datos Excel$\rightarrow$Julia$\rightarrow$Pluto via TSV.
    \item RCall.jl para pipelines mixtos Julia+R.
\end{itemize}

\section{Renderizado Quarto}

NEVEN integra Quarto 1.9.18 para renderizado de documentos reproducibles:

\begin{lstlisting}[language=bash, caption={Integración Quarto via CreateProcess}]
# Desde NEVEN Studio o celda Excel:
=NEVEN.quarto.render("reporte.qmd", "html")

# El sistema ejecuta:
quarto render reporte.qmd --to html
\end{lstlisting}

% ═══════════════════════════════════════════════════════════════════════════
\chapter{NEVEN Studio}
% ═══════════════════════════════════════════════════════════════════════════

NEVEN Studio es un modo standalone implementado como Office Add-in (HTML/JS) que extiende las capacidades de NEVEN más allá de las fórmulas de celda. El studio se comunica con un servidor HTTP Python (puerto 5555) que orquesta los motores de lenguaje.

\section{Arquitectura de NEVEN Studio}

\begin{figure}[H]
\centering
\begin{verbatim}
┌─────────────────────────────────────────────────────────────────┐
│                    NEVEN Studio (Office Add-in)                 │
├───────────┬───────────┬───────────┬───────────┬────────────────┤
│    SQL    │  Data     │   Run     │  Data     │ Presentaciones │
│   Tab     │  Studio   │  Script   │   Lab     │      Tab       │
├───────────┴───────────┴───────────┴───────────┼────────────────┤
│                      IA Tab                    │   Ayuda Tab    │
├───────────────────────────────────────────────┴────────────────┤
│                HTTP Server (Python, port 5555)                  │
├─────────────────────────────────────────────────────────────────┤
│   DuckDB   │   ControlR   │  ControlJulia  │  ControlPython    │
└─────────────────────────────────────────────────────────────────┘
\end{verbatim}
\caption{Arquitectura de NEVEN Studio}
\end{figure}

\section{Tab SQL}

Permite ejecutar consultas SQL sobre datos de Excel utilizando DuckDB embebido:

\begin{lstlisting}[language=SQL, caption={Consulta SQL en NEVEN Studio}]
-- Cargar datos desde rango de Excel
SELECT * FROM excel_range('Sheet1!A1:D100')
WHERE valor > 1000
ORDER BY fecha DESC
\end{lstlisting}

\section{Tab Data Studio}

Entorno de edición visual para transformación de datos con generación automática de código R/Python.

\section{Tab Run Script}

Editor de código con soporte para R, Julia y Python:

\begin{lstlisting}[language=R, caption={Script R en Run Script}]
# Analisis de regresion con output interactivo
library(stargazer)
modelo <- lm(y ~ x1 + x2, data = datos)
stargazer(modelo, type = "html")  # Se muestra en WebView2
\end{lstlisting}

\section{Tab Data Lab}

Pipeline de descubrimiento de datos con tres sub-tabs:
\begin{itemize}
    \item \textbf{Preguntas:} Motor de lenguaje natural para consultas sobre datos.
    \item \textbf{Explorador:} Visualización interactiva de datasets cargados.
    \item \textbf{Código:} Generación automática de código R/Python basada en preguntas.
\end{itemize}

\section{Tab Presentaciones}

Generador de presentaciones HTML autocontenidas usando reveal.js:

\begin{lstlisting}[language=javascript, caption={API de presentaciones desde Studio}]
// Crear presentacion
await NEVEN.presentation.create("Resultados Q3");

// Agregar slide con grafico Plotly
await NEVEN.presentation.addSlide("slide-1", {
    title: "Ventas por Region",
    content: plotlyJsonData,
    type: "plotly"
});

// Exportar a HTML
await NEVEN.presentation.export("presentacion.html");
\end{lstlisting}

\section{Tab IA}

Asistente de inteligencia artificial contextual con:
\begin{itemize}
    \item Integración con Claude/OpenAI via API keys configurables.
    \item Contexto automático del rango seleccionado en Excel.
    \item Generación de código R/Julia/Python basada en descripción natural.
    \item Explicación de resultados estadísticos.
\end{itemize}

\section{Tab Ayuda}

Sistema de ayuda integrado con:
\begin{itemize}
    \item Búsqueda semántica sobre funciones disponibles.
    \item Ejemplos interactivos ejecutables.
    \item Acceso al motor RAG para consultas contextuales.
\end{itemize}

% ═══════════════════════════════════════════════════════════════════════════
\chapter{Motor RAG}
% ═══════════════════════════════════════════════════════════════════════════

NEVEN incluye un motor RAG (Retrieval-Augmented Generation) para asistencia contextual y generación de código.

\section{Ontología de Dominio}

El motor RAG se sustenta en una ontología estructurada en cuatro prioridades:

\begin{table}[H]
\centering
\begin{tabular}{clc}
\toprule
\textbf{Prioridad} & \textbf{Dominio} & \textbf{Entidades} \\
\midrule
P1 & Core C++ (XLL, IPC, Memory) & 45 \\
P2 & Motores (R, Julia, Python) & 85 \\
P3 & Studio Frontend (JS/HTML) & 62 \\
P4 & Build, Deploy, Config & 38 \\
-- & Funciones estadísticas & 90+ \\
\midrule
\textbf{Total} & & \textbf{300+} \\
\bottomrule
\end{tabular}
\caption{Distribución de entidades ontológicas por prioridad}
\end{table}

\section{Arquitectura RAG}

\begin{enumerate}
    \item \textbf{Indexación:} Los archivos YAML de ontología se cargan y embeben usando sentence-transformers.
    \item \textbf{Almacenamiento:} Vectores almacenados en FAISS para búsqueda eficiente.
    \item \textbf{Recuperación:} Búsqueda semántica top-k sobre consultas del usuario.
    \item \textbf{Generación:} Contexto recuperado se inyecta en prompts para Claude/OpenAI.
\end{enumerate}

\section{Invariantes Técnicas}

La ontología define invariantes técnicas que el sistema debe respetar:

\begin{lstlisting}[language=yaml, caption={Ejemplo de invariante en ontología YAML}]
- id: INV-NX-01
  nombre: Archivos R ASCII sin BOM
  descripcion: >
    Los archivos .R DEBEN ser ASCII puro sin BOM.
    Verificar con: $bytes = [IO.File]::ReadAllBytes()
    Los primeros 3 bytes NO deben ser 239 187 191.
  severidad: critica
  verificacion: automated
\end{lstlisting}

% ═══════════════════════════════════════════════════════════════════════════
\chapter{Seguridad}
% ═══════════════════════════════════════════════════════════════════════════

\section{Sandbox de Ejecución}

El \texttt{SandboxVerifier} bloquea más de 90 patrones de ejecución peligrosa:

\begin{itemize}
    \item \textbf{R:} \texttt{system()}, \texttt{shell()}, \texttt{eval(parse())}, \texttt{do.call()}, \texttt{.Call()}, \texttt{url()}, \texttt{Sys.setenv()}
    \item \textbf{Julia:} \texttt{ccall}, \texttt{@ccall}, \texttt{cglobal}, \texttt{unsafe\_*}, \texttt{include()}, backtick literals
    \item \textbf{Python:} \texttt{exec()}, \texttt{eval()}, \texttt{subprocess}, \texttt{os.system()}, \texttt{\_\_import\_\_()}
    \item \textbf{Bypass prevention:} Detección de concatenación de strings, whitespace stripping, \texttt{getattr} patterns
\end{itemize}

\section{Seguridad del Visor WebView2}

\begin{itemize}
    \item Navegación restringida: solo \texttt{file://} dentro del directorio de datos y \texttt{about:blank}.
    \item DevTools deshabilitado por defecto (configurable para debugging).
    \item JavaScript habilitado para interactividad (Plotly, D3.js) pero sin acceso a filesystem o red.
    \item Extensión de política para \texttt{localhost:port} cuando Pluto está activo.
\end{itemize}

\section{Seguridad del Servidor HTTP}

\begin{itemize}
    \item El servidor HTTP escucha solo en localhost (127.0.0.1:5555).
    \item CORS configurado para permitir solo orígenes de Office Add-ins.
    \item Validación de inputs en todos los endpoints.
    \item Sin persistencia de credenciales (API keys en memoria, no en disco).
\end{itemize}

% ═══════════════════════════════════════════════════════════════════════════
\chapter{Correcciones Técnicas Significativas}
% ═══════════════════════════════════════════════════════════════════════════

\section{Bugs de Julia 1.12: Array I/O}

Julia 1.12 cambió el almacenamiento de arrays de tipos primitivos (\texttt{Float64}, \texttt{Int64}) a \textit{inline storage}. Las funciones \texttt{jl\_arrayset} y \texttt{jl\_array\_ptr} dejaron de funcionar correctamente para estos tipos:

\begin{lstlisting}[language=C++, caption={Fix para arrays primitivos en Julia 1.12}]
// Antes (roto en Julia 1.12):
jl_arrayset(julia_array, element, i);

// Despues (correcto):
if (array_base_type == jl_float64_type) {
    jl_array_data(julia_array, double)[i] = 
        jl_unbox_float64(element);
}
\end{lstlisting}

Adicionalmente, \texttt{jl\_array\_data} cambió de firma: \texttt{jl\_array\_data(a)} $\rightarrow$ \texttt{jl\_array\_data(a, type)}.

\section{Base.Docs.doc() en Julia 1.12}

La función \texttt{Base.Docs.doc()} lanza \texttt{MethodError} en Julia 1.12 para funciones sin docstrings formales. Se resolvió con \texttt{try/catch} en \texttt{ListFunctions()}.

\section{Python Stable ABI}

Para garantizar compatibilidad con múltiples versiones de Python sin recompilación, ControlPython.exe utiliza la Stable ABI (\texttt{python3.dll}) en lugar de la ABI versionada (\texttt{python312.dll}).

% ═══════════════════════════════════════════════════════════════════════════
\chapter{Validación y Pruebas}
% ═══════════════════════════════════════════════════════════════════════════

\section{Suite de Pruebas}

\begin{table}[H]
\centering
\begin{tabular}{lcc}
\toprule
\textbf{Categoría} & \textbf{Tests} & \textbf{Estado} \\
\midrule
Sandbox (R + Julia + Python) & 150 & 100\% \\
Property-Based Testing & 12 & 100\% \\
Confiabilidad (unit + PBT) & 45 & 100\% \\
Config, Security, Discovery & 35 & 100\% \\
Type conversions, RAII, COM & 25 & 100\% \\
Integration, Callbacks, IPC & 40 & 100\% \\
HTTP Server endpoints & 20 & 100\% \\
Misc (Result, Version) & 15 & 100\% \\
\midrule
\textbf{Total} & \textbf{342} & \textbf{100\%} \\
\bottomrule
\end{tabular}
\caption{Suite de pruebas automatizadas}
\end{table}

\section{Métricas de Rendimiento}

\begin{table}[H]
\centering
\begin{tabular}{lc}
\toprule
\textbf{Métrica} & \textbf{Valor} \\
\midrule
Latencia IPC (Named Pipes) & $<$1ms \\
Tiempo de inicio R & $\sim$2s \\
Tiempo de inicio Julia (con sysimage) & $\sim$5s \\
Tiempo de inicio Python & $\sim$1s \\
Tamaño sysimage Julia & 415MB \\
Funciones UDF registradas & 200+ \\
Entidades ontológicas RAG & 300+ \\
\bottomrule
\end{tabular}
\caption{Métricas de rendimiento del sistema}
\end{table}

\section{Evaluación Doctoral}

El sistema fue evaluado según criterios doctorales en 10 dimensiones:

\begin{table}[H]
\centering
\begin{tabular}{lcc}
\toprule
\textbf{Criterio} & \textbf{Peso} & \textbf{Nota} \\
\midrule
Arquitectura y Diseño & 15\% & 10.0 \\
Calidad del Código & 15\% & 9.5 \\
Robustez y Manejo de Errores & 10\% & 9.5 \\
Pruebas y Validación & 10\% & 10.0 \\
Documentación & 10\% & 10.0 \\
Seguridad & 10\% & 9.5 \\
Rendimiento & 5\% & 9.0 \\
Completitud Funcional & 10\% & 9.5 \\
Integración de IA/RAG & 10\% & 9.5 \\
Innovación y Contribución & 5\% & 10.0 \\
\midrule
\textbf{Nota Final Ponderada} & \textbf{100\%} & \textbf{9.71/10} \\
\bottomrule
\end{tabular}
\caption{Evaluación doctoral del sistema NEVEN}
\end{table}

\section{Evolución de Calidad}

\begin{table}[H]
\centering
\begin{tabular}{lccc}
\toprule
\textbf{Hito} & \textbf{Nota} & \textbf{Tests} & \textbf{Cambio Clave} \\
\midrule
Estado original (RJ2XCL) & 4.3 & 39 & Prototipo con deuda técnica \\
Post-correcciones & 6.8 & 119 & Seguridad, RAII, mutex \\
Post-Julia sysimage & 7.2 & 119 & Julia arranca en segundos \\
Post-Mantenibilidad & 7.9 & 165 & Constants, logging, Doxygen \\
Post-Confiabilidad & 8.1 & 200 & Health monitoring, errores \\
Post-Julia Libraries & 8.4 & 200 & 61 procedimientos Julia \\
Post-WebView2+Pluto & 8.9 & 228 & Visor, notebooks, presentaciones \\
NEVEN v3.0 (Python+Studio) & 9.2 & 290 & Python, NEVEN Studio \\
NEVEN v3.2 (RAG+DataLab) & 9.71 & 342 & RAG Engine, Data Lab, IA \\
\bottomrule
\end{tabular}
\caption{Evolución de la calificación de calidad}
\end{table}

% ═══════════════════════════════════════════════════════════════════════════
\chapter{Documentación}
% ═══════════════════════════════════════════════════════════════════════════

NEVEN incluye documentación completa de 17 capítulos en formato Docusaurus:

\begin{enumerate}
    \item Introducción y Visión General
    \item Instalación y Configuración
    \item Arquitectura del Sistema
    \item Motor R (ControlR)
    \item Motor Julia (ControlJulia)
    \item Motor Python (ControlPython)
    \item NEVEN Studio
    \item Visor WebView2
    \item Notebooks Pluto.jl
    \item Sistema de Presentaciones
    \item Motor RAG
    \item Data Lab
    \item Seguridad
    \item API Reference
    \item Troubleshooting
    \item Desarrollo y Contribución
    \item Changelog
\end{enumerate}

% ═══════════════════════════════════════════════════════════════════════════
\chapter{Conclusiones y Trabajo Futuro}
% ═══════════════════════════════════════════════════════════════════════════

\section{Conclusiones}

NEVEN demuestra que es viable construir un sistema políglota robusto que integre R, Julia y Python en Microsoft Excel mediante una arquitectura de procesos aislados. Las contribuciones principales incluyen:

\begin{enumerate}
    \item Una arquitectura IPC de cuatro capas que soporta tres lenguajes sin comprometer la estabilidad de Excel, con latencia sub-milisegundo.
    \item Más de 200 funciones UDF accesibles como fórmulas nativas de Excel.
    \item NEVEN Studio: un modo standalone con siete tabs especializados para usuarios técnicos.
    \item Un motor RAG con más de 300 entidades ontológicas para asistencia contextual.
    \item Data Lab: pipeline de descubrimiento de datos con DuckDB y lenguaje natural.
    \item Una suite de 342 pruebas automatizadas con 100\% de cobertura en paths críticos.
    \item Documentación completa de 17 capítulos en formato Docusaurus.
    \item Calificación de calidad de 9.71/10 en evaluación doctoral.
\end{enumerate}

\section{Trabajo Futuro}

\begin{itemize}
    \item \textbf{NEVEN-SIM:} Módulo de simulación Monte Carlo como XLL separado (en desarrollo).
    \item \textbf{CI/CD:} Pipeline de integración continua con tests automatizados en cada commit.
    \item \textbf{Cloud deployment:} Explorar deployment de motores en contenedores para escenarios enterprise.
    \item \textbf{GPU acceleration:} Integrar CUDA/OpenCL para cálculos intensivos en Julia.
    \item \textbf{Collaborative notebooks:} Notebooks Pluto.jl compartidos en tiempo real.
\end{itemize}

% ═══════════════════════════════════════════════════════════════════════════
\chapter*{Agradecimientos}
% ═══════════════════════════════════════════════════════════════════════════

A la Universidad de Costa Rica por el apoyo institucional durante el desarrollo de esta tesis de maestría. A BukloLAB por proporcionar el entorno de trabajo para la fase final del proyecto. Al equipo de desarrollo de BERT (Basic Excel R Toolkit) cuyo código base sirvió como punto de partida. A la comunidad de Julia por el soporte en la migración a Julia 1.12. A Anthropic y OpenAI por las APIs que potencian el asistente de IA integrado.

\end{document}
